% page 55 of the facsimile (image p-054 of the scan)
% block 162.6 x 237.7 mm ; left margin 39.4 mm
\pgc{17.780mm}{0.000mm}{
\l{\cel{52}47}
\l{}
\l{}
\l{\sou{ateliero}. Loko en qua laboras piktisto, skultisto, fotografisto,}
\l{\cel{2}+artizano. - DF.}
\l{}
\l{\sou{atencar}. (\sou{trans.}) Direktar sua tota penso ad.- DEFIS.}
\l{}
\l{\sou{atentar}. (\sou{trans}.) Krimino-probar. - DEFIS.}
\l{}
\l{\sou{ateromo}. (\sou{patol}.) Degenero di la tuniko interna di la arterii,}
\l{\cel{3}qua implikas lokale la existo di amasi de kolesterino e di}
\l{\cel{3}kalko-sali. - DEFIRS.}
\l{}
\l{\sou{atestar}. (\sou{trans}.)Deklarar ke ulo esas exakta, reala. - DEFIRS.}
\l{}
\l{\sou{atiko}. Supera parto di edifico, ordinare ornita per pilastri,}
\l{\cel{3}e di qua la rolo esas disimular la tekto. - EF.}
\l{}
\l{\sou{atingar}.\cel{2}(\sou{trans.})\cel{2}Sucesar tushar to quo distas. - EFIS.}
\l{}
\l{\sou{atipa}. (\sou{biol}.) Qua esas plu o min fora de la tipo normala.}
\l{\cel{3}- DEFIS.}
\l{}
\l{\sou{atlanto}. (\sou{arkitekt.)}\cel{2}Statuo di viro, nuda e drapirita, qua}
\l{\cel{3}subtenas entablamento o kornico. - DEFIRS.}
\l{}
\l{\sou{atlaso}. Kolekturo de mapi, de imaji anatomiala, anexita a libro.}
\l{\cel{3}- DEFIRS.}
\l{}
\l{\sou{atleto}.\cel{2}La homo qua, en la ludi publika, kombatis pri luktado,}
\l{\cel{3}e pri boxado (en Grekia antiqua). - DEFIRS.}
\l{}
\l{\sou{atmosfero}. Aero-zono qua cirkondas la Terglobo. - DEFIRS.}
\l{}
\l{\sou{atolo}. Insulo ek koralio. - DFS.}
\l{}
\l{\sou{atomo}. Ta parteto di korpo simpla, quan onu konsideris kom}
\l{\cel{3}nepartigebla, e qua esas la elemento-quanto maxim mikra}
\l{\cel{3}apta divenar parto de kombinuro. - DEFIRS.}
\l{}
\l{\sou{atonio}. Ne-existo di vivozeso. - DEFIRS.}
\l{}
\l{\sou{- atr -}. Sufixo qua indikas simileso.}
\l{}
\l{\sou{atraktar}.\cel{2}(\sou{trans.}) Kaptar per ruzo-stroko. - DEFS.}
\l{}
\l{\sou{atribuar}.\cel{2}(\sou{trans., ad)}\cel{2}Donar (ulo) ad ulu kom lua parto.}
\l{\cel{3}- DEFIS.}
\l{}
\l{\sou{atributo}. (\sou{gram.}) Termino di la propoziciono, qua expresas}
\l{\cel{3}la eso-maniero quan onu afirmas pri la subjekto. - DEFIRS.}
\l{}
\l{\sou{atricar}. (\sou{netrans)}. (\sou{teol.}) Repentar pri sua peko sive pro}
\l{\cel{3}shamo, sive pro timo a puniseso. - DEFIS.}
}
